% !TeX root = ../教材题库_全部条目_转写工作稿.tex
% Source: XJ-B1-C06-RV-P267-Y0346
% 第6章 / 章末复习 复习题六 / 教材 p.267 / PDF 第274页
% OCR未校对：公式与文字均可能有误；未完成数学TeX转写；禁止直接发给学生
% 原题及相邻原解截图（核对用）：../题目原图/XJ-B1-C06-RV-P267-Y0346-01.webp
% 本文件只在转写完成后作为题干使用；OCR 条目尚不是可靠题干。
\begin{bankitem}{复习题六 第13题}{第6章 / 章末复习 复习题六 / 教材 p.267 / PDF 第274页}{XJ-B1-C06-RV-P267-Y0346}{OCR未校对：公式与文字均可能有误；未完成数学TeX转写；禁止直接发给学生}
\begin{ocrdraft}
[教材p.267 / PDF第274页]
\par
13．为了鼓励居民节约用水，市政府计划调整居民生活用水收费方案，拟确
\par
定一个合理的月用水量标准：一位居民的月用水量不超过的部分按平价收费，
\par
超出xt的部分按议价收费．为了解居民用水情况，通过抽样，获得了100位居民
\par
每人的月均用水量（单位：t)，将数据按照[0，0．5），
\par
[4，4．5）分
\par
成9组，制成了如图所示的频率分布直方图．
\par
组距
\par
0．52
\par
住16
\par
0．12
\par
0．08
\par
0．04
\par
0
\par
511．522．533．544．5
\par
月均用水量/t
\par
（第13题）
\par
（D求直方图中a的值，并估计居民月均用水量的中位数；
\par
（2）设该市有30万居民，估计全市居民中月均用水量不低于3t的人数，并说
\par
明理由；
\par
（3）若该市政府希望使85％的居民月用水量不超过标准t,估计的值，并
\par
说明理由．
\end{ocrdraft}
\end{bankitem}
